\section{Zero-shot 与 Few-shot In-context Learning}

\subsection{GPT 系列的演进}

GPT 系列是后训练研究的重要背景。让我们回顾其演进：

\begin{table}[H]
\centering
\begin{tabular}{lccc}
\toprule
\textbf{模型} & \textbf{参数量} & \textbf{训练数据} & \textbf{关键突破} \\
\midrule
GPT-1 (2018) & 117M & 4.6 GB & 预训练+微调范式 \\
GPT-2 (2019) & 1.5B & 40 GB & Zero-shot 能力涌现 \\
GPT-3 (2020) & 175B & 600 GB & Few-shot in-context learning \\
\bottomrule
\end{tabular}
\caption{GPT 系列模型的演进}
\end{table}

\begin{figure}[H]
\centering
\includegraphics[width=0.95\textwidth]{slides-images/slide-008.jpg}
\caption{GPT-1 架构：Decoder-only Transformer，12 层，Next Token Prediction 损失\protect\footnotemark}
\end{figure}
\footnotetext{Slides 第9页。}

\subsection{Zero-shot Learning}

Zero-shot learning 指的是模型在\textbf{没有看到任何任务示例}的情况下，直接完成任务。这需要我们巧妙地将任务表述为文本补全问题。

\begin{figure}[H]
\centering
\includegraphics[width=0.95\textwidth]{slides-images/slide-012.jpg}
\caption{Zero-shot 学习示例：通过巧妙的 prompt 设计让模型完成 QA 和推理任务\protect\footnotemark}
\end{figure}
\footnotetext{Slides 第13页。}

核心思路：既然模型会补全文本，我们就把任务\textbf{伪装成}文本补全。例如：
\begin{itemize}
    \item \textbf{问答}：将上下文和问题拼接起来，让模型补全答案
    \item \textbf{摘要}：在文章末尾加上 ``TL;DR''，模型自然会生成摘要
    \item \textbf{消歧}：将两个候选替换分别填入句子，比较模型赋予的对数概率
\end{itemize}

\begin{figure}[H]
\centering
\includegraphics[width=0.95\textwidth]{slides-images/slide-016.jpg}
\caption{Zero-shot 摘要：在文章后添加 ``TL;DR''，模型自动生成摘要\protect\footnotemark}
\end{figure}
\footnotetext{Slides 第17页。}

\begin{importantbox}{GPT-2 的 Zero-shot 能力}
GPT-2 无需任何任务特定的微调，仅靠预训练就能在多个基准任务上达到有竞争力的表现。这是一个标志性的结果——意味着足够大的预训练模型已经``内化''了大量任务的解法。
\end{importantbox}

\subsection{Few-shot In-context Learning}

GPT-3 带来了更令人惊叹的能力：\textbf{Few-shot in-context learning}。只需要在 prompt 中提供几个示例，模型就能``学会''新任务——完全\textbf{不需要梯度更新}。

\begin{figure}[H]
\centering
\includegraphics[width=0.95\textwidth]{slides-images/slide-019.jpg}
\caption{Few-shot learning：在上下文中放入几个翻译示例，模型即可执行翻译任务\protect\footnotemark}
\end{figure}
\footnotetext{Slides 第20页。}

\begin{figure}[H]
\centering
\includegraphics[width=0.95\textwidth]{slides-images/slide-020.jpg}
\caption{GPT-3 的 Few-shot 性能：随着示例数量增加，性能逐步逼近甚至超越专用模型\protect\footnotemark}
\end{figure}
\footnotetext{Slides 第21页。}

\begin{knowledgebox}{In-context Learning 与传统 Fine-tuning 的区别}
传统的 fine-tuning 需要：(1) 收集大量标注数据，(2) 对模型进行多轮梯度更新，(3) 可能损害其他任务的性能。而 in-context learning 仅需在推理时的输入中放入几个示例，模型参数完全不变。这种能力随模型规模增大而\textbf{涌现}——小模型几乎没有这种能力，但 175B 参数的 GPT-3 表现出了惊人的 few-shot 性能。
\end{knowledgebox}

\subsection{Chain-of-Thought 提示}

即使有了 few-shot 能力，模型在复杂推理任务上仍然表现不佳。\textbf{Chain-of-Thought (CoT)} 提示技术通过让模型``展示推理过程''来大幅提升性能。

\begin{figure}[H]
\centering
\includegraphics[width=0.95\textwidth]{slides-images/slide-023.jpg}
\caption{Chain-of-Thought 提示：左侧为标准 few-shot（直接给答案），右侧为 CoT（展示推理步骤）\protect\footnotemark}
\end{figure}
\footnotetext{Slides 第24页。}

核心思想有两种形式：

\textbf{1. Few-shot CoT}：在示例中不仅给出答案，还给出推理过程。模型在回答新问题时也会模仿这种推理链。

\textbf{2. Zero-shot CoT}：无需任何示例，仅在 prompt 中添加 ``Let's think step by step''，模型就会自动展开推理过程。

\begin{figure}[H]
\centering
\includegraphics[width=0.95\textwidth]{slides-images/slide-025.jpg}
\caption{Zero-shot CoT：仅添加``Let's think step by step''就能大幅提升推理准确率\protect\footnotemark}
\end{figure}
\footnotetext{Slides 第26页。}

\begin{figure}[H]
\centering
\includegraphics[width=0.95\textwidth]{slides-images/slide-026.jpg}
\caption{CoT 的性能对比：Zero-shot 仅 17.7\%，Zero-shot CoT 达 78.7\%，Few-shot CoT 接近超越 supervised SOTA\protect\footnotemark}
\end{figure}
\footnotetext{Slides 第27页。}

\begin{warningbox}{提示工程的本质}
提示工程的本质是\textbf{理解预训练数据的分布}。当你设计一个 prompt 时，你实际上是在寻找：互联网上什么样的文本之后会跟着我想要的那种回答？ ``Let's think step by step'' 之所以有效，是因为互联网上以这种方式开头的文本通常后面跟着详细的推理过程。然而，这种``欺骗模型''的方式本质上是不令人满意的——我们不应该需要技巧来让模型做正确的事。
\end{warningbox}

\subsection{In-context Learning 的局限性}

尽管提示技术取得了巨大成功，但存在根本性的局限：

\begin{itemize}
    \item \textbf{上下文窗口有限}：能放入的示例数量受限（虽然现代模型的上下文越来越长）
    \item \textbf{需要``技巧''}：用户需要精心设计 prompt 才能获得好的结果，这对非专业用户不友好
    \item \textbf{无法处理复杂任务}：多步推理、创意写作等任务难以仅通过提示解决
    \item \textbf{模型不是助手}：预训练模型的目标是预测文本，而非帮助用户
\end{itemize}

\begin{figure}[H]
\centering
\includegraphics[width=0.95\textwidth]{slides-images/slide-028.jpg}
\caption{预训练模型的限制：问``向6岁小孩解释登月''，GPT-3 却继续生成更多问题而非回答\protect\footnotemark}
\end{figure}
\footnotetext{Slides 第29页。}

这个例子清楚地展示了问题所在：预训练模型不是在\textbf{回答}你的问题，而是在\textbf{补全}一段看起来合理的文本。这种``不与用户意图对齐''的问题，正是后续方法要解决的核心挑战。

\subsection{本章小结}

\begin{itemize}
    \item Zero-shot 和 Few-shot 能力随模型规模增大而涌现，是预训练规模化的重要成果
    \item Chain-of-Thought 提示通过引导模型``思考''来提升推理能力
    \item 但提示技术本质上是在``欺骗''模型，用户体验差且能力有限
    \item 我们需要更根本的方法让模型从``文本补全器''变为``用户助手''
\end{itemize}

%% ========== Part 3: 指令微调 ==========

